\begin{singlespace}
\section*{Resumen}

\textbf{Título:} \tit\footnote{Trabajo de grado}

\textbf{Autores:} \nam\footnote{\fac.\ \esc.\\ Director: \tdir\ \dir, Codirector: \tcdir\ \cdir }

\textbf{Palabras clave:} { Computación cuántica, Simulación cuántica, Gestión de memoria,\\ Compresión de datos }

\textbf{Descripción:}

La simulación clásica de circuitos cuánticos sigue siendo un recurso indispensable para diseñar, validar y analizar algoritmos cuánticos. No obstante, en los simuladores de estado completo su alcance práctico queda limitado por el crecimiento exponencial de la memoria requerida para almacenar el vector de estado. En respuesta a esta restricción, el presente trabajo estudia la compresión sin pérdida como mecanismo de gestión de memoria en simuladores cuánticos tipo Schr\"odinger, con el objetivo de reducir la huella de almacenamiento sin introducir alteraciones numéricas en la simulación.

La propuesta comprende la selección de un simulador base y de una biblioteca de compresión adecuada, el diseño de una arquitectura de almacenamiento por bloques con descompresión selectiva y caché LRU, su implementación en TMFQSfullstate mediante Blosc2 y su evaluación experimental frente a una ejecución de referencia sin compresión y a una estrategia con pérdida controlada basada en ZFP. La validación se realizó con circuitos representativos de comportamientos contrastantes de compresibilidad, en particular Grover y la transformada cuántica de Fourier, para instancias entre 20 y 25 qubits.

Los resultados muestran que la estrategia sin pérdida basada en Blosc logra reducciones significativas de memoria en escenarios favorables, especialmente en Grover, manteniendo coincidencia exacta con la referencia no comprimida y error absoluto máximo nulo en todos los casos evaluados. También evidencian que la eficacia de la compresión depende de la estructura del estado cuántico y del patrón de acceso al vector, por lo que no puede asumirse un beneficio uniforme entre algoritmos. En conjunto, el trabajo demuestra que la compresión sin pérdida es una alternativa viable para ampliar la capacidad práctica de simulación cuando la memoria es la restricción dominante y el estado conserva regularidades aprovechables.
\end{singlespace}

\newpage

\begin{singlespace}
\section*{Abstract}

\textbf{Title:} {Optimization of memory management in quantum simulators through data\\ compression without loss of precision}\footnote{Degree Thesis}

\textbf{Authors:} \nam\footnote{Faculty of Physics-Mechanics Engineering. School of Systems Engineering and Computer
Science.\\ Advisor: \tdir\ \dir, Co-Advisor: \tcdir\ \cdir }

\textbf{Key words:} { Quantum Computing, Quantum Simulation, Memory Management, Data Compression }

\textbf{Description:}

Classical quantum-circuit simulation remains an essential resource for designing, validating, and analyzing quantum algorithms. In full-state simulators, however, practical scale is limited by the exponential memory required to store the state vector. In response to that restriction, this thesis studies lossless compression as a memory-management mechanism for Schr\"odinger-style quantum simulators, aiming to reduce storage demands without introducing numerical alterations into the simulation.

The proposed approach includes the selection of a base simulator and a suitable compression library, the design of a block-based storage architecture with selective decompression and LRU caching, its implementation in TMFQSfullstate using Blosc2, and its experimental evaluation against both an uncompressed reference and a controlled-loss strategy based on ZFP. Validation was conducted with representative circuits exhibiting contrasting compressibility behavior, namely Grover's algorithm and the quantum Fourier transform, for instances between 20 and 25 qubits.

The results show that the Blosc-based lossless strategy achieves significant memory reductions in favorable scenarios, especially for Grover, while preserving exact agreement with the uncompressed reference and zero maximum absolute error in all the cases. They also show that compression effectiveness depends on the structure of the quantum state and on state-vector access patterns, so benefits cannot be assumed to be uniform across algorithms. Overall, the thesis demonstrates that lossless compression is a viable alternative for extending practical simulation capacity when memory is the dominant constraint and the state preserves exploitable regularity.
\end{singlespace}

\newpage
